\section{Conclusions}
The conservative FDM correctly represents the divergence-form diffusion
operator when adjacent rows share each positive face conductance and
Dirichlet data are eliminated symmetrically. Its discrete energy identity
establishes positive definiteness, and its M-matrix sign pattern gives a
maximum principle. The smooth manufactured benchmark, independently
checked source term, four-level error table, and algebraic tolerance sweep
support approximately second-order accuracy for $a=1+x$.

The layered test separates algebraic structure from physical accuracy.
Harmonic averaging reproduces the piecewise-linear exact solution and
constant flux at the nodes and faces to round-off. Arithmetic averaging
retains symmetry, local discrete conservation, and the maximum principle,
yet its nodal error is first order and its maximum interface-flux error
stays near $0.662$ for the tested two-dimensional boundary data. A single
flux-jump statistic would have hidden this defect.

The matched FEM comparison gives the same approximately second-order order
on the smooth problem, with different measured constants and costs. The
fitted layered FEM is also exact for this particular solution. Contrast
increases the unscaled FDM condition number and unpreconditioned CG work;
Jacobi substantially reduces the contrast effect in the tested sequence.

The experiments are limited to a Cartesian square, a smooth affine
coefficient, and a straight two-material interface. Curved interfaces,
non-affine smooth coefficients, and general heterogeneous patterns require
separate verification. The supplied code and raw data permit the reported
tables and figures to be regenerated. Individual contribution statements
must be supplied by the team members from their own records.
